\section{Local polynomials and the two asymptotic questions}\label{cs:introduction}
A Casimir source is a homogeneous polynomial obtained from a finite
matrix operator associated with a pair of strands. A diagram combines
these local polynomials by identifying their face variables and extracting
one coefficient. There are two different analytic questions. First, how
large can a local polynomial be over all face phases and the full real
spectral interval? Second, does the shared coefficient have a nonzero
leading term, and can one bound the rest of its spectral domain?

We answer the first question for each of the three degree sectors of the
local polynomial. For the second question we select the central sector
at every matched edge of a triangulated sphere and assume that the
associated triangle variational problem has stationary positive radii.
Theorem~\ref{lc:main} gives the local upper bounds;
Theorem~\ref{sa:main} gives a holomorphic Gaussian asymptotic for the
shared central coefficient and a strict spectral gap away from zero.
Their hypotheses and conclusions are kept separate. In particular, the
second theorem does not estimate the complementary degree sectors.

The proofs start with exact finite cofactor identities. Signed transfer
matrices and logarithmic integrals give the uniform local bounds.
For the nonzero asymptotic, Euler--Maclaurin summation with Cauchy
estimates constructs local analytic actions, and a contour deformation
selects their contributing Gaussian. Positive finite matrices then
control the other face phases before a common coefficient is extracted.
The one-variable summation and saddle formulas are recorded in
\cite[\S2.10(i), \S2.4(iv)]{DLMF}; the uniform errors, contour
contribution, and shared multivariable Gaussian are proved below.

The variables of a local face polynomial have a prescribed order:
\(x,y\) are the outer variables, \(b\) the input middle variable, and
\(c\) the output middle variable. Variables shared by different boxes
must later be identified. This paper estimates one actual polynomial
before that shared extraction. The estimates hold for all phases of
the four variables, rather than only at their positive real values.

Fix an integer \(N\ge3\), and put
\[
 d=N-1,\quad m=\lfloor N/2\rfloor,\quad J=\lfloor d/2\rfloor,\quad
 \theta=\pi/N,\quad q=e^{i\theta},\quad Q=q^2.
\]
The finite Fourier theorem in \cite[Proposition 4.1 and
equations (4.4)--(4.5)]{FinitePaper} identifies the physical pair Casimir
\(\Omega\), one \(N\)-dimensional residue block at a time, with the
polynomial operator specified in Section~\ref{lc:source}. The identification
uses a finite Fourier transform, a transpose and a diagonal change of basis;
it is stated precisely below. It also proves \(\chi_N(\Omega)=0\), where
\(\chi_N(z)=2T_N(z/2)+2\) and \(T_N(\cos t)=\cos Nt\).
The full resolvent numerator is the polynomial
\[
 V_N(\sigma;\Omega)
       =\frac{\chi_N(\sigma)-\chi_N(\Omega)}{\sigma-\Omega}.
\]
The fraction is polynomial division, not an inverse at a spectral root.
Its degree in \(\sigma\) is \(d\); it retains the first derivatives at
double roots as well as the values at distinct roots.

Write \(H_V(x,b,c,y;\sigma)\) for its homogeneous face polynomial.
Section~\ref{lc:source} gives a division-free definition and fixes its
sign. A decreasing path from row \(j\) to row \(i\) contributes the
middle monomial \(b^ic^{d-j}\). When \(N=2m\), its central sector
has \(i\le m-1,j\ge m\); when \(N=2m+1\), it has \(i\le m,j\ge m\).
The left sector has \(j<m\); the right consists of the remaining
paths above the central cut. Denote these three polynomials by
\(H_{V,\mathrm C},H_{V,\mathrm L},H_{V,\mathrm R}\).
This is a partition of the full numerator paths, not a new spectral
operator.

Set
\[
 \mathcal B(t)=-\int_0^t\log|2\sin v|\,dv,\qquad
 e(a)=\frac{\mathcal B(\pi a)+\mathcal B(\pi/2-\pi a)}{\pi},
 \quad 0\le a\le\tfrac12.
\]
For nonnegative \(w_1,w_2,w_3\), not all zero, define
\begin{equation}\label{lc:entropy}
 \mathcal Q(w_1,w_2,w_3)=
 \max_{\substack{a_i\ge0,\ \sum_i a_i=1/2\\a_i=0\ {\rm if}\ w_i=0}}
 \left\{\sum_i e(a_i)+\sum_{w_i>0}a_i\log w_i\right\}.
\end{equation}\aliascurrentequation{\label{sa:entropy}}
A variable fixed at zero has exactly zero load. It is never replaced
silently by a small positive radius. The logarithmic integrals are
finite at their endpoints. To fix their normalization, put
\(I=\int_0^{\pi/2}\log\sin t\,dt\). Reflection and the double-angle
formula give \(2I=I-(\pi/2)\log2\). Hence
\(\mathcal B(0)=\mathcal B(\pi/2)=0\), and \(e(0)=e(1/2)=0\).

In the theorem the lowercase letters denote nonnegative moduli.
For \(x+y>0\), put
\begin{equation}\label{lc:rates}
\begin{aligned}
 Q_b&=\mathcal Q(x,b,y),&Q_c&=\mathcal Q(x,c,y),&X&=Q_b+Q_c,\\
 A_{\mathrm L}&=\max\{\log c,Q_b+\tfrac12\log c\}&&(c>0),\\
 A_{\mathrm R}&=\max\{\log b,Q_c+\tfrac12\log b\}&&(b>0).
\end{aligned}
\end{equation}
If \(c=0\), the left polynomial is zero and its rate is interpreted
as \(-\infty\); the same convention applies on the right when \(b=0\).

\begin{theorem}[Uniform local sector bounds]\label{lc:main}
Fix nonnegative face moduli \((x,b,c,y)\) with \(x+y>0\). For
\(s=\mathrm C,\mathrm L,\mathrm R\),
\begin{equation}\label{lc:main-rate}
 \limsup_{N\to\infty}\frac1N
 \log\sup_{\substack{\sigma\in[-2,2]\\
 |X_0|=x,\ |B_0|=b,\ |C_0|=c,\ |Y_0|=y}}
 |H_{V,s}(X_0,B_0,C_0,Y_0;\sigma)|
 \le
 \begin{cases}
 X&s=\mathrm C,\\
 A_{\mathrm L}&s=\mathrm L,\\
 A_{\mathrm R}&s=\mathrm R.
 \end{cases}
\end{equation}
The limit superior is over all integer levels, including both parities
and composite levels. The statement is uniform when the zero pattern
is fixed and every nonzero modulus varies in a fixed compact subset
of \((0,\infty)\). Thus, for every \(\varepsilon>0\), there is one
threshold \(N_0\) on such a compact set for which the indicated rate
plus \(\varepsilon\) bounds the normalized logarithm at every \(N\ge N_0\).
\end{theorem}

The supremum precedes the limit, so sampled spectral zeros are included.
The proof never divides a finite numerator by such a zero, and never
replaces a negative tail coefficient by its absolute value inside the
recurrence.

The admissible vertex with load \(1/2\) on \(b\) gives
\begin{equation}\label{lc:spectator}
 \mathcal Q(x,b,y)\ge\tfrac12\log b\quad(b>0),\qquad
 \max\{X,A_{\mathrm L},A_{\mathrm R}\}=\max\{X,\log b,\log c\},
\end{equation}
with undefined logarithms omitted. Each mixed expression in
\eqref{lc:rates} is at most \(X\). This bounds the full local polynomial. A shared coefficient must still
be estimated at one common set of face radii, because the same variable
can occur in several local factors.

